% Part: counterfactuals
% Chapter: minimal-change-semantics
% Section: transitivity

\documentclass[../../../include/open-logic-section]{subfiles}

\begin{document}

\olfileid{cnt}{min}{cpo}

\olsection{ప్రతివిపర్యయం}

భౌతిక, కఠిన సోపాధికాలు
తమ ప్రతివిపర్యయాలకు తుల్యాలు.
ప్రతివాస్తవాలు అలా కావు.
క్రాట్జర్ ఇచ్చిన ఉదాహరణ ఇది:
\begin{quote}
  గోథె 1832లో మరణించి
  ఉండకపోయినా, ఇప్పుడు
  మరణించి ఉండేవారు.

  గోథె ఇప్పుడు మరణించి
  ఉండకపోతే, 1832లో
  మరణించి ఉండేవారు.
\end{quote}
మొదటి వాక్యం సత్యం: మనుషులు
వందల సంవత్సరాలు జీవించరు.
రెండోది స్పష్టంగా అసత్యం:
గోథె ఇప్పుడు మరణించి ఉండకపోతే
ఇంకా సజీవంగానే ఉండేవారు;
అందువల్ల 1832లో మరణించి
ఉండలేరు.

\begin{ex}\ollabel{ex:contraposition-counterex}
  గోళ అర్థవిచారంలో ప్రతివిపర్యయ
  నిగమనం చెల్లదు; అంటే $p
  \cif q \Entails/ \lnot q \cif \lnot p$.
  $p$ను ``గోథె 1832లో
  మరణించలేదు'' అని, $q$ను
  ``గోథె ఇప్పుడు మరణించి
  ఉన్నారు'' అని భావించండి.
  దీనిని $W =
  \{w, w_1, w_2\}$,
  $O_w = \{\{w\}, \{w, w_1\},
  \{w, w_1, w_2\}\}$,
  $V(p) = \{w_1, w_2\}$,
  $V(q) = \{w, w_1\}$
  గల నమూనా $\mModel{M_1}
  = \tuple{W, O, V}$లో
  చూపవచ్చు.
  \sourcecorrection{OLTECNTCPO-001}{మూల నమూనా వివరణలో ఒక్క లోకం చుట్టూ గోళ వ్యవస్థకు $O$ అని వ్రాసింది; నిర్వచనంలో $O$ ప్రతి లోకానికి గోళ వ్యవస్థను ఇచ్చే ప్రమేయం, ఆ లోకం $w$ వద్ద గోళ వ్యవస్థ $O_w$. ఇక్కడ $O_w$గా సరిచేశాం.}
  కాబట్టి $w$ వాస్తవ లోకం:
  అందులో గోథె 1832లో మరణించి,
  ఇప్పటికీ మరణించి ఉన్నారు;
  $w_1$ సమీప లోకం: అందులో
  గోథె, ఉదాహరణకు 1833లో
  మరణించి, ఇప్పటికీ మరణించి
  ఉన్నారు; $w_2$ దూరపు లోకం:
  అందులో గోథె ఇంకా సజీవం.
\begin{figure}
\begin{center}
\begin{tikzpicture}[scale=.6,layerwidth=1.5]\tiny
  \spheresystem{3}
  \begin{scope}
    \clip \propositionplot{0}{.8}{50}{4.5} -- (4.5,-4.5) -- (-4.5,-4.5) -- (-4.5,4.5) -- (4.5,4.5);
    \spherelayer{3}
  \end{scope}
  \propositionintersect{0}{2}{110}{5.5}
  \proposition{0}{.8}{50}{4.5}
  \path (0,0) node[world] {$w$};
  \spherepos{0}{2}{node[world] {$w_1$}}
  \spherepos{-50}{3}{node[world] {$w_2$}}
  \spherepos{28}{3}{node {$q$}}
  \spherepos{42}{3}{node {$\lnot q$}}
  \spherepos{-55}{4}{node {$p$}}
  \spherepos{-67}{4}{node {$\lnot p$}}
\end{tikzpicture}
\caption{ప్రతివిపర్యయానికి ప్రతిదృష్టాంతం}
\ollabel{fig:contraposition}
\end{center}
\end{figure}
$p$ సత్యమయ్యే లోకాన్ని
కలిగిన గోళం $S =
\{w, w_1\}$ ఉంది;
దాని అన్ని లోకాల వద్ద
$p \lif q$ సత్యం.
కాబట్టి $\mSat{M}{p
  \cif q}[w]$. అయితే
$\lnot q$ సత్యమయ్యే
లోకాన్ని కలిగిన గోళం
$\{w, w_1,
w_2\}$లో ఒక లోకం,
అంటే~$w_2$, ఉంది.
అక్కడ $q$ అసత్యం,
$p$ సత్యం; కాబట్టి
$\mSat/{M}{\lnot q \lif \lnot p}[w_2]$.
\end{ex}

\end{document}
